\paragraph*{04.05.06.03 Erreichung der Qualitätsziele des Projekts verifizieren und erforderliche korrektive
und/oder präventive Maßnahmen empfehlen}
\kcilabel{para:04.05.06.03_ErreichungDerQualitätsziele}{04.05.06.03}
\label{para:04.05.06.03_ErreichungDerQualitätsziele}

% Task
\par Nach der erfolgreichen Entwicklung, Implementierung und Entwicklertestung der \gls{ricefw}-Objekte im \gls{epa} und \gls{zp1} in Welle 1 war es von entscheidender Bedeutung, die Verifizierung der Erreichung der definierten Qualitätsziele durch die abnehmenden Stakeholder sicherzustellen. Darüber hinaus sollten erforderliche korrigierende und/oder präventive Maßnahmen identifiziert und empfohlen werden, um eine kontinuierliche Verbesserung der Qualität der \gls{ricefw}-Entwicklung zu gewährleisten.

% Action
\par Zur Sicherstellung der Qualität entwickelte ich gemeinsam mit dem Programmtestmanager einen umfassenden Qualitätsmanagementplan, der die spezifischen Qualitätsanforderungen des \gls{wsr} adressierte.\endnote{Im \gls{s4} agiere ich als stv. \gls{tm}} Der Plan umfasste die Definition verbindlicher Qualitätsstandards, die Festlegung geeigneter Qualitätskontroll- und Qualitätssicherungsprozesse sowie die Implementierung von Maßnahmen zur kontinuierlichen Verbesserung (vgl. \autoref{fig:TestpläneSST} und \autoref{fig:UATDatenDoku}). Nichtsdestotrotz wurden erforderlichen Tests im \gls{fb} unterlassen (vgl. \autoref{fig:R2HStatusSSTJuni2025})

% Result
\par Trotz Wie in \autoref{para:04.04.05.02_OwnershipUndCommitment} dargestellt, wurde der Qualitätsmanagementplan erfolgreich umgesetzt und trug maßgeblich zu einer exzellenten Qualität der \gls{ricefw}-Entwicklung im \gls{epa} und \gls{zp1} in Welle~1 bei, wie der Fehlerstand bei \gls{hypc} Ende zeigt (vgl. \autoref{tab:statusuebersicht}).

% Reflect
\par  Die strukturierte Umsetzung des Qualitätsmanagementplans führte zu einer signifikanten Steigerung der Produktqualität sowie zu einer erhöhten Zufriedenheit der beteiligten Stakeholder.